\chapter{Théorème de confluence HMT–MD : forme électrogravitationnelle, système d’Einstein–Schrödinger et coordination de l’action, de la matière, des champs, de la géométrie et de l’information}
\label{chap:parte-iv-51}

L’unification s’exprime comme compatibilité verticale d’expressions distinctes d’un même état : couplage et vide, action et gravité, géométrie et information, écran et torsion, horloge et horizon. Le théorème conserve les domaines et les applications de chaque réponse.

\section{Théorème de confluence des réponses dimensionnelles}
\label{rm:ch:confluencia-general}

\subsection{Antécédent commun et produit fibré}

Les résultats du volume proviennent de lecteurs hétérogènes appliqués à un même état généalogique enrichi. Soit

\begin{equation}
 X=\bigl(X_{\APP},X_{\TRIT},X_{\TPK},
          \mathcal C,\mathcal B,\mathcal R,\mathcal M,\partial X\bigr)
 \label{rm:eq:confluence-state}
\end{equation}

l’état qui conserve les cylindres, la retenue, le parcours, la mémoire et la frontière. Sur \(X\) sont construits les lecteurs typés

\begin{equation}
 P_i:X\longrightarrow Y_i,
 \qquad
 i\in I:=\{\alpha,{\rm act},{\rm ang},108,\square,{\rm obs},{\rm Per}\}.
 \label{rm:eq:confluence-objects}
\end{equation}

Leurs codomaines représentent respectivement le générateur dodécaphasique de couplage, le couple d’action, l’opérateur angulaire bicouche, l’horloge \(108\), l’écran cubique, l’observateur interne et le mode de Perron. Pour chaque lecteur, on considère son graphe, muni de la projection canonique

\begin{equation}
 \Gamma_i:=\{(x,y_i)\in X\times Y_i:P_i(x)=y_i\},
 \qquad
 \pi_i:\Gamma_i\longrightarrow X,
 \quad (x,y_i)\longmapsto x.
 \label{rm:eq:confluence-reader-graphs}
\end{equation}

L’objet maître est le produit fibré bien typé de ces graphes :

\begin{equation}
 \boxed{
 \mathfrak M_X=
 \mathop{\times}_{X,\,i\in I}\Gamma_i
 =\left\{\bigl(x,(y_i)_{i\in I}\bigr):
 P_i(x)=y_i\ \text{pour tout }i\in I\right\}.}
 \label{rm:eq:master-fiber-product}
\end{equation}

Les projections \(\pi_i\) fixent matériellement le même antécédent. Un élément de \(\mathfrak M_X\) est une famille de sorties compatibles accompagnées de l’état qui les engendre.

\begin{definition}[Publication dimensionnelle]
Une publication dimensionnelle est une composition
\begin{equation}
 \mathfrak M_X\xrightarrow{\;P\;}\mathcal L
 \xrightarrow{\;\operatorname{ev}_{\mathcal U}\;}\RR,
 \label{rm:eq:dimensional-publication}
\end{equation}
où \(\mathcal L\) est une droite dimensionnelle typée et \(\operatorname{ev}_{\mathcal U}\) est l’évaluation finale dans une carte d’unités. L’application \(P\) construit d’abord l’invariant et sa généalogie.
\end{definition}

\subsection{Indépendance horizontale et compatibilité verticale}

Deux lecteurs issus d’un même antécédent peuvent satisfaire des relations croisées sans se factoriser l’un par l’autre. Cette distinction régit les doubles feuillets du vide, les réalisations positive et orientée, la double projection archimédienne et exceptionnelle, ainsi que les lectures géométrique, informationnelle et de saveur de la transition \(6\to8\).

\begin{theorem}[Confluence sans factorisation horizontale]
Soient \(P_i:\mathfrak M_X\to Y_i\), \(i=1,\ldots,r\), des publications construites sur le même état. Supposons que chaque \(P_i\) conserve ses données typées et que les relations \(R_j(P_1,\ldots,P_r)=0\) soient démontrées après la construction des lecteurs. Alors le morphisme conjoint
\begin{equation}
 P=(P_1,\ldots,P_r):\mathfrak M_X\longrightarrow
 Y_1\times\cdots\times Y_r
 \label{rm:eq:joint-publication}
\end{equation}
est verticalement compatible avec \(X\). Une factorisation \(P_i=F_{ij}\circ P_j\) exige en outre que \(P_j\) sépare toutes les fibres que \(P_i\) distingue. La seule existence de \(R_j=0\) n’implique pas cette séparation.
\end{theorem}

\begin{proof}
La compatibilité verticale résulte de la propriété universelle du produit fibré des graphes \(\Gamma_i\) : toutes les composantes ont la même image dans \(X\). Si \(P_i=F_{ij}\circ P_j\), alors \(P_j(x)=P_j(x')\) impose \(P_i(x)=P_i(x')\). Par conséquent, toute paire \(x,x'\) ayant la même lecture \(j\) et des lectures \(i\) différentes réfute la factorisation. Une relation scalaire \(R_j=0\) ne fait que restreindre l’image conjointe ; elle ne garantit pas l’inclusion des noyaux nécessaire à la factorisation.
\end{proof}

Le théorème explique pourquoi le produit des feuillets peut fixer le cône causal tandis que leur quotient conserve l’orientation, et pourquoi l’aire totale peut coexister avec un Hamiltonien modulaire qui distingue la provenance \(90/120\).

\subsection{1re confluence : couplage, vide et secteur électrofaible}

Soit

\begin{equation}
 D=D_A=\det T,
 \qquad
 Q(D)=-\frac9{25}\log D=4\pi\alpha_{\HMT}.
 \label{rm:eq:confluence-q-d}
\end{equation}

La mémoire orientée s’exprime par

\begin{equation}
 \rho=e^{2C^*_{\rm rad}},
 \qquad
 q_-=\sqrt{D\rho},
 \qquad
 q_+=\sqrt{D/\rho}.
 \label{rm:eq:confluence-qpm}
\end{equation}

Le calcul fonctionnel

\begin{equation}
 F(z)=\frac{1-z^{90}}{1-z^{120}},
 \qquad
 r_\pm=F(q_\pm)
 \label{rm:eq:confluence-F}
\end{equation}

produit

\begin{equation}
 \widehat\mu=r_-^2,
 \quad
 \widehat\varepsilon=r_+^2,
 \quad
 \widehat Z=\frac{r_-}{r_+},
 \quad
 \widehat c=(r_-r_+)^{-1}.
 \label{rm:eq:confluence-vacuum}
\end{equation}

À l’interface électrofaible au niveau des arbres,

\begin{equation}
 g^2=\frac{Q(D)}{1-D},
 \qquad
 g'^2=\frac{Q(D)}{D},
 \label{rm:eq:confluence-gauge-couplings}
\end{equation}

d'où

\begin{equation}
 \boxed{
 (1-D)g^2=Dg'^2=Q(D),
 \qquad
 \frac1{g^2}+\frac1{g'^2}=\frac1{Q(D)}.}
 \label{rm:eq:confluence-ew-master}
\end{equation}

La première égalité exprime le même couplage dans deux orientations de jauge ; la seconde le retrouve comme somme de compliances. Le déterminant \(D\) régit le mélange pair et \(\rho\) conserve l’orientation impaire.

La coordonnée de Higgs intègre cette dernière :

\begin{equation}
 \lambda_H=\frac{Q(D)}{8(1-D)}D^{-3}\rho^{5/3}.
 \label{rm:eq:confluence-higgs}
\end{equation}

En définissant

\begin{equation}
 \Xi=\frac{8(1-D)D^3\lambda_H}{Q(D)},
 \label{rm:eq:confluence-xi}
\end{equation}

est reconstruit

\begin{equation}
 \rho=\Xi^{3/5},
 \qquad
 \widehat Z=
 \frac{F(\sqrt D\,\Xi^{3/10})}
      {F(\sqrt D\,\Xi^{-3/10})}.
 \label{rm:eq:confluence-vacuum-higgs}
\end{equation}

L’équation \cref{rm:eq:confluence-vacuum-higgs} constitue un contrôle croisé entre orientation scalaire et impédance du vide. Sa portée provient de la construction parallèle des deux lecteurs sur \(X\).

\subsection{2e confluence : action, gravité et forme électrogravitationnelle}

Soit

\begin{equation}
 Q_P=E_P,
 \qquad
 \ell_P,
 \qquad
 Q_P\ell_P=\hbar c,
 \qquad
 \frac{\ell_P}{Q_P}=\frac{G}{c^4}.
 \label{rm:eq:confluence-planck-pair}
\end{equation}

Le produit est le quantum d’échange \(\hbar c\) ; le quotient est la compliance gravitationnelle. Introduisons

\begin{equation}
 \mathscr T_P=\frac{Q_P}{\ell_P}=\frac{c^4}{G},
 \qquad
 \mathscr C_P=\frac{\ell_P}{Q_P}=\frac{G}{c^4}.
 \label{rm:eq:confluence-tension-compliance}
\end{equation}

L’équation d’Einstein adopte la forme constitutive

\begin{equation}
 G_{\mu\nu}+\Lambda g_{\mu\nu}
 =8\pi\mathscr C_P T_{\mu\nu}.
 \label{rm:eq:confluence-einstein-constitutive}
\end{equation}

Pour une charge \(q\), la coordonnée électrique portée au plan de l’énergie est

\begin{equation}
 \mathcal Q(q)=Q_P\sqrt{\frac{Z}{4\pi\hbar}}\,q.
 \label{rm:eq:confluence-electric-coordinate}
\end{equation}

Newton et Coulomb se réunissent alors dans la forme bilinéaire

\begin{equation}
 \boxed{
 V(r)=\frac{\mathscr C_P}{r}
 \bigl(\mathcal Q_1\mathcal Q_2-E_1E_2\bigr).}
 \label{rm:eq:confluence-newton-coulomb}
\end{equation}

Le cône \(\mathcal Q^2=E^2\) reproduit la condition statique d’extrémalité dans la carte d’Einstein--Maxwell. Le produit des feuillets du vide régit le cône causal ; leur quotient régit le relèvement électrique et l’orientation.

L’identité croisée avec le secteur électrofaible est

\begin{equation}
 \boxed{
 \frac{\hbar^2}{8\pi\ell_P^2}
 \frac{(\kappa/\mu)}{e^2}
 =\frac1{4\pi\alpha_{\HMT}}
 =\frac1{g^2}+\frac1{g'^2}.}
 \label{rm:eq:confluence-einstein-maxwell-ew}
\end{equation}

Cette égalité conserve sa valeur sous l’inversion du feuillet d’action : \(G\) et \(e^2\) se transforment de manière contragrédiente, tandis que \(G\hbar\), \(Ge^2\) et \(\ell_P^2\) restent invariants.

\subsection{3e confluence : géométrie, information et mélange}

Soit \(N=N_{90}+N_{120}\) le nombre total d’excitations de bord et

\begin{equation}
 D_{90/120}=90N_{90}+120N_{120}.
 \label{rm:eq:confluence-modular-provenance}
\end{equation}

L’aire lit \(N\), tandis que le Hamiltonien holographique modulaire lit \(D_{90/120}\). Comme

\begin{equation}
 \det\begin{pmatrix}1&1\\90&120\end{pmatrix}=30,
 \label{rm:eq:confluence-modular-determinant}
\end{equation}

le couple reconstruit exactement

\begin{equation}
 N_{90}=4N-\frac{D_{90/120}}{30},
 \qquad
 N_{120}=\frac{D_{90/120}}{30}-3N.
 \label{rm:eq:confluence-sector-reconstruction}
\end{equation}

La géométrie totale et l’énergie modulaire pondérée conservent donc la quantité et la provenance.

L’état angulaire bicouche

\begin{equation}
 \rho_y=\frac{e^{-yR}}{2\cosh y}
 \label{rm:eq:confluence-rhoy}
\end{equation}

possède une entropie paire et une divergence orientée

\begin{equation}
 S(\rho_y)=\log(2\cosh y)-y\tanh y,
 \qquad
 D(\rho_y\Vert\rho_{-y})=2y\tanh y.
 \label{rm:eq:confluence-info-orientation}
\end{equation}

L’involution \(y\mapsto-y\) est représentée dans le secteur de saveur par

\begin{equation}
 R_{\rm CKM}
 =M_{\rm CKM}\operatorname{diag}(1,-1,1)M_{\rm CKM}^{-1},
 \qquad
 R_{\rm CKM}^2=I,
 \qquad
 \det R_{\rm CKM}=-1.
 \label{rm:eq:confluence-rckm}
\end{equation}

Barbero lit la composante impaire de la transition \(6\to8\), l’entropie lit sa distribution paire et \(R_{\rm CKM}\) réalise la même inversion dans le codomaine de mélange. Ce sont des représentations coordonnées d’une charnière commune, dotées de domaines et d’observables propres.

\subsection{4e confluence : écran, torsion et mémoire hélicoïdale}

Dans le quotient lorentzien \(Q_{\square}\), la soudure et la réponse positive satisfont

\begin{equation}
 \boxed{
 (\beta_m^{\square})^*\Lambda_m^{\square}\beta_m^{\square}
 =\frac{56}{81}I_{Q_{\square}}.}
 \label{rm:eq:confluence-screen-energy}
\end{equation}

Dans le registre profond, l’opérateur de vis vérifie

\begin{equation}
 \boxed{
 V^{-1}\mathscr SV=(2+\sqrt3)\ii\,\mathscr S.}
 \label{rm:eq:confluence-screw}
\end{equation}

La première identité mesure le coût minimal de réponse de l’écran ; la seconde transporte l’orientation et la profondeur. Catalan apparaît comme le deuxième moment de la composante quaternaire. Si \(N_{\rm dep}\delta_n=n\delta_n\) désigne l’opérateur de profondeur du registre hélicoïdal, alors

\begin{equation}
 G_{\rm Cat}=\Tr(N_{\rm dep}^{-2}\mathcal Q_{\chi,-4}).
 \label{rm:eq:confluence-catalan}
\end{equation}

La loi de Perron

\begin{equation}
 M_n=\left(\frac{a_n}{a_0}\right)^{-6}
 \label{rm:eq:confluence-sixth-law}
\end{equation}

fournit la covariance exacte d’une densité quadratique de spin. Le réalisateur de Cartan doit transporter \cref{rm:eq:confluence-screen-energy,rm:eq:confluence-screw} vers un courant et une contorsion dont le mode homogène satisfait \cref{rm:eq:confluence-sixth-law}. Cet enchaînement situe l’interface physique dans une flèche unique et falsifiable.

\subsection{Confluence entre périodicité de l’horloge, énergie d’horizon et coût informationnel}

L’horloge \(108\) fixe

\begin{equation}
 E_P=k_B\Theta_{\rm clk}\log3.
 \label{rm:eq:confluence-planck-trit}
\end{equation}

Pour la région de Sitter de rayon \(R=\ell_P\), l’énergie du vide, la température et l’entropie satisfont

\begin{equation}
 E_\Lambda=T_H S_H=\frac{E_P}{2}
 =\frac{k_B\Theta_{\rm clk}\log3}{2}.
 \label{rm:eq:confluence-horizon-half}
\end{equation}

Il en résulte

\begin{equation}
 E_\Lambda t_P=\frac{\hbar}{2},
 \qquad
 E_\Lambda T_{108}=\frac{h}{2},
 \qquad
 \frac{S_H}{k_B}=\pi.
 \label{rm:eq:confluence-half-action}
\end{equation}

Le volume du vide, l’aire de l’horizon, la température du cycle et le coût informationnel d’un trit expriment la même énergie. La multiplicité \(e^{S_H/k_B}=e^\pi\) coïncide exactement avec la valeur propre expansive de \(e^{\pi J_\varphi}\) ; promouvoir cette égalité en équivalence de microétats exige un entrelaceur préservant la mesure et la dynamique.

\subsection{Confluence fibrée des réponses HMT–MD : vide, action, gravité, information, mélange, mémoire et cosmologie}

\begin{theorem}[Confluence des réponses HMT--MD]
L’état enrichi \(X\) et les primitives déclarées dans chaque lecteur étant fixés, les publications de vide, d’action, de gravité, d’information, de mélange, d’écran, de mémoire hélicoïdale et de cosmologie sont les projections ou compositions définies sur le produit fibré de graphes \cref{rm:eq:master-fiber-product}. Leurs relations croisées sont engendrées par \cref{rm:eq:confluence-ew-master,rm:eq:confluence-newton-coulomb,%
rm:eq:confluence-einstein-maxwell-ew,rm:eq:confluence-sector-reconstruction,%
rm:eq:confluence-info-orientation,rm:eq:confluence-screen-energy,%
rm:eq:confluence-screw,rm:eq:confluence-horizon-half}.
Chaque évaluation numérique est terminale et chaque mémoire orientée demeure dans le facteur qui l’engendre.
\end{theorem}

\begin{proof}
La construction de \(\mathfrak M_X\) fixe l’antécédent commun. Les équations du vide résultent du calcul fonctionnel de \(T\) ; les équations électrofaibles résultent de \(D\) et \(Q(D)\) ; la forme électrogravitationnelle résulte des deux identités du couple de Planck ; la reconstruction \(90/120\) résulte de l’inversion d’une matrice entière de déterminant \(30\) ; l’orientation informationnelle résulte du logarithme de \(\rho_y\) ; l’énergie d’écran résulte de la réponse unique dans \(Q_{\square}\) ; la vis résulte des deux covariances de Weyl et de Pell ; et l’égalité d’horizon résulte des formules géométriques de Sitter spécialisées en \(R=\ell_P\). Aucune preuve n’utilise la décimale terminale pour sélectionner la branche. Le théorème d’indépendance horizontale interdit de transformer les relations croisées en factorisations non démontrées.
\end{proof}

\subsection{Corollaire de confluence métrologique entre les réalisations HMT et SI}

Toute sortie admet le quintuplet indissociable

\begin{equation}
 \mathcal K(P)=
 \bigl(
 \text{antécédent},
 \text{équation},
 \text{invariant},
 \text{signification},
 \text{valeur terminale}
 \bigr).
 \label{rm:eq:confluence-quintuple}
\end{equation}

L’atlas métrologique du volume énumère ces quintuplets. Son contenu probatoire est le suivant : l’antécédent fixe la généalogie ; l’équation fixe l’opération ; l’invariant démontre ce qui subsiste ; la signification explicite la lecture physique ; la valeur terminale permet la reconnaissance externe. Séparer une colonne détruit l’information causale ; conserver les cinq permet d’insérer le résultat dans une carte d’unités sans faire de cette carte un fondement.

\begin{proofstatusbox}
Le produit fibré, les identités algébriques et le théorème d’indépendance horizontale constituent une formalisation exacte. Les équations internes citées conservent le statut démontré dans les chapitres qui contiennent leurs démonstrations complètes. Les lectures d’Einstein--Maxwell, de Sitter, d’Einstein--Cartan et électrofaible demeurent dans leurs interfaces déclarées lorsqu’elles nécessitent un réalisateur supplémentaire. Le théorème organise ces couches sans les promouvoir par similitude numérique.
\end{proofstatusbox}

\begin{falsifierbox}
La confluence est réfutée si deux composantes attribuées au même élément de \(\mathfrak M_X\) proviennent en réalité d’états incompatibles, si l’une des identités génératrices échoue, si une évaluation terminale intervient dans la sélection de son propre antécédent, ou si une mémoire distinguée par un autre lecteur est perdue. Une prétendue factorisation horizontale est réfutée par toute paire d’états ayant le même lecteur intermédiaire et des sorties finales différentes. Les promotions physiques sont en outre confrontées à leurs propres réalisateurs, identités de jauge, de Bianchi, de raffinement et conditions aux limites.
\end{falsifierbox}


\section{Réalisation physique discrète par produit fibré et entrelacement de Cartan}

Le produit fibré commun, ses composantes positives et orientées, la localisation réversible et le critère de promotion lisse fournissent l’entrelaceur matériel entre la confluence abstraite et ses réalisations physiques. Le développement est conservé intégralement, avec ses obstructions explicites.

\begingroup
\let\chapter\section
\let\section\subsection
\let\subsection\subsubsection
\let\subsubsection\paragraph
\let\clearpage\relax
\let\cleardoublepage\relax
\section[Réalisation physique discrète sur les voies TPK]{Réalisation physique
discrète : lectures positive et orientée sur les voies TPK, domaine commun et
localisation réversible}
\label{sec:int68-realizacion-fisica-discreta}

La forme spectrale, la réalisation de Hilbert, la section dimensionnelle et le
transport de Cartan ne sont pas des décorations ajoutées à une sortie classique. Ce sont
quatre représentations coordonnées d'une même voie TPK enrichie. Pour
exprimer cette unité causale, il faut d'abord construire leur domaine commun ; il ne
suffit pas d'écrire quatre flèches ayant le même symbole de départ.

\subsection{Domaine commun comme produit fibré}

Soit \(\mathsf P_{\rm TPK}^{\rm enr}\) la catégorie des états TPK préparés
et des voies enregistrées compatibles avec leurs extrémités, leur frontière et leurs restrictions.
Soient
\[
 \mathsf D_{\rm sp}:=\mathsf{SurvPath}_{\rm HMT},\qquad
 \mathsf D_Q:=\operatorname{Dom}(\mathcal Q),\qquad
 \mathsf D_D:=\operatorname{Dom}(\mathfrak D_{\rm HMT}),\qquad
 \mathsf D_C:=\operatorname{Dom}(\rho_C^{\rm disc}),
\]
avec leurs foncteurs d'oubli fidèles
\[
 U_i:\mathsf D_i\longrightarrow\mathsf P_{\rm TPK}^{\rm enr}
 \qquad(i\in\{{\rm sp},Q,D,C\}).
\]
Chaque \(U_i\) n'oublie que la réalisation et conserve l'état et la voie
TPK sous-jacents. Nous définissons
\begin{equation}
 \mathsf P_{\rm CT}^{+}
 :=
 \mathsf D_{\rm sp}
 \times_{\mathsf P_{\rm TPK}^{\rm enr}}\mathsf D_Q
 \times_{\mathsf P_{\rm TPK}^{\rm enr}}\mathsf D_D
 \times_{\mathsf P_{\rm TPK}^{\rm enr}}\mathsf D_C.
 \label{eq:int68-pct-pullback}
\end{equation}
Un objet est donc un état préparé muni des quatre
réalisations sur le même antécédent ; une flèche est une voie enregistrée
dont l'image sous-jacente coïncide dans les quatre facteurs. Si les
présentations ne sont identifiées qu'à isomorphisme près, on utilise le produit
fibré bicatégorique correspondant ; dans la présentation fixée du registre,
on adopte le produit fibré strict de
\eqref{eq:int68-pct-pullback}.

\begin{proposition}[Existence d'un objet dans le domaine commun]
\label{prop:int68-pct-no-vacio}
La catégorie \(\mathsf P_{\rm CT}^{+}\) n'est pas vide. Il n'est pas nécessaire de
particulariser la voie électronique centrale pour le prouver : toute section
coinductive certifiée, équipée de sa flèche identité et de ses données neutres, admet
des relèvements compatibles aux quatre domaines de
\eqref{eq:int68-pct-pullback}.
\end{proposition}

\begin{proof}
Le \cref{thm:generacion-monodromica-conjunta-rev7} fournit, par exemple,
une section survivante préparée \(x_\pi\). Son identité
\(\operatorname{id}_{x_\pi}\) est simultanément : une voie survivante ;
un morphisme de \(\mathsf P_{\rm TPK}\), que \(\mathcal Q\) envoie sur
l'identité bornée de \(\mathcal Q(x_\pi)\) ; la voie vide de signature nulle,
que le \cref{thm:funtor-dimensional-minimo} envoie sur
\((0,0,0,1)\) ; et l'identité du groupoïde, que le
\cref{thm:realizacion-discreta-cartan-rev6} envoie sur l'élément neutre
\((I,0)\) de Poincaré. Par oubli, ces quatre relèvements donnent le même objet
\(x_\pi\) et le même morphisme \(\operatorname{id}_{x_\pi}\). Leur quadruplet
définit, par la propriété du produit fibré, un objet de
\(\mathsf P_{\rm CT}^{+}\).
\end{proof}

\begin{definition}[Adaptateurs au domaine commun]
\label{def:int68-pct-positive}
Les quatre adaptateurs sont les projections canoniques
\[
 \begin{aligned}
 F_{\rm sp}:\mathsf P_{\rm CT}^{+}&\longrightarrow
     \mathsf{SurvPath}_{\rm HMT},\\
 F_Q:\mathsf P_{\rm CT}^{+}&\longrightarrow\operatorname{Dom}(\mathcal Q),\\
 F_D:\mathsf P_{\rm CT}^{+}&\longrightarrow
     \operatorname{Dom}(\mathfrak D_{\rm HMT}),\\
 F_C:\mathsf P_{\rm CT}^{+}&\longrightarrow
     \operatorname{Dom}(\rho_C^{\rm disc}).
 \end{aligned}
\]
En particulier, \(F_{\rm sp}\) conserve la famille compatible des
restrictions finies de la section et de la voie ; il ne reconstruit pas cette famille
à partir de son spectre.
\end{definition}

\begin{lemma}[Fonctorialité des adaptateurs]
\label{lem:int68-adaptadores}
Les quatre adaptateurs préservent les identités et la composition.
\end{lemma}

\begin{proof}
L'identité d'un objet du produit fibré est le quadruplet des identités
sur le même état TPK ; chaque projection renvoie l'identité du facteur.
La composition est définie composante par composante entre quadruplets qui
se projettent sur la même concaténation de voies. Projeter avant ou après
composition donne, par définition, le même résultat.
\end{proof}

\subsection{\(4\) composantes et produit positif}

Les théorèmes de réalisation spectrale, hilbertienne, dimensionnelle et de Cartan
fournissent
\[
 \mathfrak S:\mathsf{SurvPath}_{\rm HMT}
      \longrightarrow\mathbf{SpecRoute},
 \qquad
 \mathcal Q:\operatorname{Dom}(\mathcal Q)
      \longrightarrow\mathsf{Hilb}_{\mathbb R,J},
\]
\[
 \mathfrak D_{\rm HMT}:\operatorname{Dom}(\mathfrak D_{\rm HMT})
      \longrightarrow\mathbf B\mathfrak M_D^+,
 \qquad
 \rho_C^{\rm disc}:\operatorname{Dom}(\rho_C^{\rm disc})
      \longrightarrow\mathbf B G_C,
\]
où
\[
 G_C:=\operatorname{Spin}^{+}(1,3)\ltimes\mathbb R^{1,3}
\]
et
\begin{equation}
 \mathfrak M_D^+
 :=\mathbb N\times\mathbb R_{\geq0}^{2}\times\mathbb R_{>0},
 \qquad
 (n,a,b,\lambda)\odot(n',a',b',\lambda')
 =(n+n',a+a',b+b',\lambda\lambda').
 \label{eq:int68-monoide-dimensional-positivo}
\end{equation}
Ici \(\mathbf BM\) est la catégorie à un objet associée au monoïde
\(M\). Le premier foncteur est celui du
\cref{thm:functor-realizacion-espectral-numero} ; le deuxième est la
réalisation de Hilbert définie dans le chapitre des postulats ; le troisième est
le \cref{thm:funtor-dimensional-minimo} ; et le quatrième est le
\cref{thm:realizacion-discreta-cartan-rev6}.

Nous définissons les restrictions communes
\[
 \mathfrak S_+:=\mathfrak S\circ F_{\rm sp},\qquad
 \mathcal Q_+:=\mathcal Q\circ F_Q,\qquad
 \mathfrak D_+:=\mathfrak D_{\rm HMT}\circ F_D,\qquad
 \rho_{C,+}^{\rm disc}:=\rho_C^{\rm disc}\circ F_C.
\]

\begin{theorem}[Réalisation discrète positive sur le domaine commun]
\label{thm:int68-producto-positivo}
L'application
\begin{equation}
 \mathfrak R_{\rm MD}^{\rm disc,+}
 :=(\mathfrak S_+,\mathcal Q_+,\mathfrak D_+,
       \rho_{C,+}^{\rm disc})
 \label{eq:int68-producto-positivo}
\end{equation}
est un foncteur
\[
 \mathsf P_{\rm CT}^{+}\longrightarrow
 \mathbf{SpecRoute}\times\mathsf{Hilb}_{\mathbb R,J}
 \times\mathbf B\mathfrak M_D^+\times\mathbf BG_C.
\]
On n'affirme pas que toute voie TPK appartienne à \(\mathsf P_{\rm CT}^{+}\).
\end{theorem}

\begin{proof}
Par le \cref{lem:int68-adaptadores}, chaque projection \(F_i\) est un foncteur ;
par les théorèmes cités, il en va de même de
\(\mathfrak S,\mathcal Q,\mathfrak D_{\rm HMT}\) et de
\(\rho_C^{\rm disc}\) dans leurs domaines déclarés. Leurs compositions
préservent donc les identités et la concaténation. La composition dans le codomaine
produit se calcule composante par composante, ce qui prouve
\eqref{eq:int68-producto-positivo}.
\end{proof}

La composante de Hilbert a pour cible générale
\(\mathsf{Hilb}_{\mathbb R,J}\), dont les morphismes sont des opérateurs réels
bornés qui entrelacent les structures complexes. Soit
\(\mathsf P_{\rm CT}^{u}\subseteq\mathsf P_{\rm CT}^{+}\) la sous-catégorie
des flèches \(\gamma:x\to y\) pour lesquelles
\begin{equation}
 \mathcal Q_+(\gamma)^*\mathcal Q_+(\gamma)=I_{\mathcal Q_+(x)},
 \qquad
 \mathcal Q_+(\gamma)\mathcal Q_+(\gamma)^*=I_{\mathcal Q_+(y)}.
 \label{eq:int68-unitariedad-doble}
\end{equation}
Les identités et les produits d'opérateurs unitaires sont encore
unitaires ; cette condition définit donc une sous-catégorie, et ce n'est qu'en elle que
l'on obtient
\begin{equation}
 \mathcal U_{\rm reg}:=
 \mathcal Q_+\big|_{\mathsf P_{\rm CT}^{u}}:
 \mathsf P_{\rm CT}^{u}\longrightarrow
 \mathsf{UHilb}_{\mathbb R,J}.
 \label{eq:int68-restriccion-unitaria}
\end{equation}
L'inversibilité unilatérale ou la conservation de la norme sur un sous-ensemble ne
remplacent pas les deux identités de \eqref{eq:int68-unitariedad-doble}.

La notation à sept entrées se décompose en quatre composantes
typées : \(Q\) est \(\mathcal Q_+\) ;
\(U\) est la restriction \eqref{eq:int68-restriccion-unitaria} ; \(A\) est la
coordonnée d'action de \(\mathfrak D_+\) ; \(C\) est la conjugaison déclarée
dans le sous-domaine où elle entrelace la composante hilbertienne et l'orientation ;
\(\mathcal E\) et \(\mathcal B\) sont des applications de lecture constitutives subordonnées aux
données dimensionnelles ; et \(N_{\rm sp}\) est une propriété quadratique de la
réalisation scindée. Aucune des quatre dernières n'est ajoutée au produit comme
un cinquième foncteur sans domaine ni codomaine.

\subsection{Lectures positive et orientée}

Soit \(e\mapsto e^\vee\) le renversement des arêtes du graphe dirigé
sous-jacent. Dans \(\mathsf P_{\rm CT}^{+}\), \(e^\vee\) est une autre arête et n'est
pas encore déclarée inverse catégorique de \(e\). Un poids
\(w(e)=w(e^\vee)>0\) étant fixé, nous définissons sur une voie
\(\gamma=e_1\cdots e_r\)
\[
 V_+(\gamma):=\sum_{j=1}^r w(e_j),
 \qquad
 V_{\rm or}(\gamma):=
 \sum_{j=1}^r\varepsilon(e_j)w(\underline e_j),
\]
où \(\underline e_j\) est l'arête non orientée et
\(\varepsilon(e^\vee)=-\varepsilon(e)\). La première application de lecture conserve
la variation totale positive ; la seconde, le transport net orienté. Toutes deux sont
additives par concaténation et reçoivent la même voie TPK.

\begin{proposition}[Non-factorisation du lecteur positif par le lecteur orienté]
\label{prop:int68-no-factorizacion}
Supposons que l'identité et le lacet
\(ee^\vee:=e^\vee\circ e\) appartiennent au domaine construit. Il n'existe pas de
fonction \(H\) sur \(\operatorname{im}V_{\rm or}\) telle que
\(V_+=H\circ V_{\rm or}\).
\end{proposition}

\begin{proof}
On a
\[
 V_{\rm or}(\operatorname{id})=0
 =V_{\rm or}(ee^\vee),
 \qquad
 V_+(\operatorname{id})=0,
 \qquad
 V_+(ee^\vee)=2w(e)>0.
\]
Une même entrée \(0\) obligerait \(H(0)\) à prendre deux valeurs distinctes.
\end{proof}

Ce résultat est une non-factorisation concrète dans le domaine construit ; ce
n'est pas un théorème universel sur tout opérateur d'entrelacement enrichi possible.
En particulier, il n'autorise pas à supprimer une application de lecture. L'unité réside dans
l'antécédent TPK commun et non dans une compression horizontale entre ses sorties.

\subsection{Localisation réversible et réalisation orientée}

Soit \(\mathsf P_{{\rm CT},{\rm rev}}^+\) la sous-catégorie engendrée par les
arêtes pour lesquelles les quatre conditions suivantes ont été vérifiées :
\begin{enumerate}
 \item \(\mathfrak S_+(e)\) est un isomorphisme de
       \(\mathbf{SpecRoute}\) et
       \(\mathfrak S_+(e^\vee)=\mathfrak S_+(e)^{-1}\) ;
 \item \(\mathcal Q_+(e)\) satisfait
       \eqref{eq:int68-unitariedad-doble} et
       \(\mathcal Q_+(e^\vee)=\mathcal Q_+(e)^*\) ;
 \item la donnée dimensionnelle orientée provient d'une cochaîne additive avec
       \(d_{\rm or}(e^\vee)=-d_{\rm or}(e)\) et prend ses valeurs dans
       \(\operatorname{Gr}(\mathfrak M_D^+)\) ;
 \item les cocycles de Cartan sont normalisés, additifs et antisymétriques
       sous renversement.
\end{enumerate}
Nous formons le groupoïde
\[
 \Pi_\vee(\mathsf P_{{\rm CT},{\rm rev}}^+)
\]
obtenu à partir de la complétion en groupoïde en imposant en outre
\([e^\vee]=[e]^{-1}\). La composante dimensionnelle orientée prend ses valeurs dans
\begin{equation}
 \operatorname{Gr}(\mathfrak M_D^+)
 \cong\mathbb Z\times\mathbb R^2\times\mathbb R_{>0},
 \label{eq:int68-grothendieck-dimensional}
\end{equation}
avec pour inverse
\((n,a,b,\lambda)^{-1}=(-n,-a,-b,\lambda^{-1})\). Ce lecteur signé n'est pas
la grandeur positive \(V_+\) rebaptisée : ses valeurs sur les arêtes
renversées sont distinctes.

Pour la composante de Cartan, soient \(c_g,c_s\) des cocycles comme au point 4,
soit \(R_4\in\operatorname{Spin}^{+}(1,3)\), soit \(R_4u=u\) et soit
\(\ell_0>0\). Alors
\[
 \rho_C^{\rm disc}(\gamma)
 =\bigl(R_4^{c_g(\gamma)},\ell_0c_s(\gamma)u\bigr).
\]
La loi
\((R,a)(S,b)=(RS,a+Rb)\), l'égalité \(R_4u=u\) et l'additivité des
cocycles donnent
\[
 \rho_C^{\rm disc}(\gamma_2\gamma_1)
 =\rho_C^{\rm disc}(\gamma_2)\rho_C^{\rm disc}(\gamma_1),
 \qquad
 \rho_C^{\rm disc}(\gamma^\vee)
 =\rho_C^{\rm disc}(\gamma)^{-1}.
\]

\begin{theorem}[Réalisation discrète orientée]
\label{thm:int68-realizacion-orientada}
Sous les quatre conditions précédentes, les composantes descendent de manière
unique en un foncteur
\[
 \mathfrak R_{\rm MD}^{\rm disc,or}:
 \Pi_\vee(\mathsf P_{{\rm CT},{\rm rev}}^+)
 \longrightarrow
 \operatorname{Core}(\mathbf{SpecRoute})
 \times\mathsf{UHilb}_{\mathbb R,J}
 \times\mathbf B\operatorname{Gr}(\mathfrak M_D^+)
 \times\mathbf BG_C.
\]
En particulier,
\[
 \mathfrak R_{\rm MD}^{\rm disc,or}(\gamma^{-1})
 =\mathfrak R_{\rm MD}^{\rm disc,or}(\gamma)^{-1}.
\]
\end{theorem}

\begin{proof}
Les conditions 1 et 2 rendent inversibles les images spectrale et
hilbertienne ; la condition 3 prend ses valeurs dans un groupe ; et la composante de Cartan
prend déjà ses valeurs dans \(G_C\). Les quatre respectent la relation génératrice
\([e^\vee]=[e]^{-1}\). Par la propriété universelle de la localisation en
groupoïde, chacune descend de manière unique ; leur produit est le foncteur
indiqué. L'identité d'inversion se vérifie composante par composante.
\end{proof}

L'application de lecture positive \(\mathfrak R_{\rm MD}^{\rm disc,+}\) et l'orientée
\(\mathfrak R_{\rm MD}^{\rm disc,or}\) sont sœurs : aucune n'est définie comme
compression de l'autre. Le passage ultérieur d'une valeur orientée à son module n'est
pas, en général, monoïdal et ne coïncide avec une lecture positive que dans un cône
sans annulation. De même, cette paire ne remplace pas la double projection
fondatrice : l'évaluation archimédienne et la branche
Paley--Witt--Golay--Leech--\(V^\natural\)--Monstre--Moonshine demeurent
des applications de lecture sœurs du même état enrichi. Les quatre composantes
construites dans cette section ne remplacent aucune de ces deux projections.

\subsection{Entrelacement de Cartan et limite lisse}

Soit
\[
 \mathcal A_{\ell_0}
 =\begin{pmatrix}\omega&\ell_0^{-1}e\\0&0\end{pmatrix}
\]
une connexion de Cartan déclarée. L'équation
\begin{equation}
 \operatorname{Hol}_{\mathcal A_{\ell_0}}
   \bigl(\mathfrak R_C(\gamma)\bigr)
 =\rho_C\bigl(\operatorname{Hol}_{\rm TPK}(\gamma)\bigr)
 \label{eq:int68-cuadrado-cartan}
\end{equation}
est exacte dans le modèle discret lorsque les deux holonomies sont définies par le
même produit ordonné d'arêtes.

\begin{criterion}[Promotion à une réalisation lisse]
\label{crit:int68-cartan-suave}
L'équation \eqref{eq:int68-cuadrado-cartan} ne définit un prolongement
lisse que si l'on prouve en outre l'indépendance vis-à-vis de la subdivision, la convergence du
produit ordonné et la récupération locale d'une cotétrade et d'une connexion avec
les domaines et la régularité déclarés.
\end{criterion}

Le terme « physique » désigne ici le fait que les quatre composantes publient
le spectre, la réalisation de l'état, les dimensions et le transport de Cartan sans
perdre leur antécédent TPK\@. Il n'affirme ni unicité, ni équivalence de catégories,
ni couverture de toute voie, ni carte SI automatique, ni théorie des champs complète.

\begin{criterion}[Obstructions aux réalisations positive, unitaire et orientée]
\label{crit:int68-obstrucciones-realizacion}
La réalisation positive n'est pas définie si l'un des quatre
adaptateurs manque, si leurs images ne possèdent pas la même voie TPK sous-jacente ou si une
composante ne préserve pas l'identité et la composition. La restriction unitaire
n'existe pas si l'une quelconque des deux égalités de
\eqref{eq:int68-unitariedad-doble} manque. La réalisation orientée ne descend pas au
groupoïde si une image spectrale n'est pas isomorphe, si une relation localisée
change d'image, si une grandeur positive est traitée comme un inverse, si un
renversement ne change pas le signe des données orientées ou si le cocycle de
Cartan cesse d'être normalisé, additif ou antisymétrique. Il n'existe pas de
prolongement lisse compatible si le résultat dépend de la subdivision ou si
la limite ne retrouve pas la cotétrade et la connexion. Toute utilisation d'une coordonnée
externe pour sélectionner la voie viole le domaine construit.
\end{criterion}

\paragraph{Domaine, codomaine et statut des équations de confluence HMT–MD}
Le produit positif est démontré sur le produit fibré commun et la
réalisation orientée, sur la localisation réversible sous les quatre
conditions explicites. Étendre le domaine à d'autres voies exige de construire leurs
quatre relèvements sur un même antécédent. Le prolongement lisse exige,
en outre, les propriétés de subdivision, de convergence et de récupération locale du
\cref{crit:int68-cartan-suave} ; une vérification finie ne remplace pas ces
propriétés.

\endgroup

\paragraph{Passage de la réalisation de Cartan–Holst à la confluence horizon–aire–information par conservation de l’action} La première réalisation constitutive de cette confluence est le vide électromagnétique--électrofaible, où les fonctions symétrique, antisymétrique et déterminantale expriment propagation et couplages.
